\documentclass[10pt,a4paper]{amsart}
\usepackage[margin=19mm]{geometry}
\usepackage{amsmath,amssymb,mathtools}
\usepackage[scr=rsfs,cal=euler]{mathalfa}
\usepackage{xfrac}
\usepackage[unicode,hidelinks,bookmarks=false]{hyperref}
\newcommand{\ie}{i.e.\ }
\newcommand{\cf}{cf.\ }
\newcommand{\resp}{respectively\ }
\newcommand{\Subsection}{\subsection}
\providecommand{\nobreakdash}{\nobreak\hbox{-}}
\setlength{\emergencystretch}{2em}
\multlinegap0pt
\usepackage{amssymb}
\usepackage{mathtools}
\newtheorem{theorem}{Theorem}
\newtheorem{lemma}{Lemma}
\DeclareMathOperator{\III}{III}
\title{A Diophantine Criterion for the Shafarevich--Tate Groups of Elliptic Curves from Heron Triangles}
\author{Vinodkumar Ghale}
\date{Source: arXiv:2301.03486v5 (CC BY 4.0)}
\begin{document}
\maketitle

\noindent\emph{Source and licence.} A Diophantine Criterion for the Shafarevich--Tate Groups of Elliptic Curves from Heron Triangles. Version arXiv:2301.03486v5 (CC BY 4.0), distributed by arXiv under the Creative Commons Attribution 4.0 International licence, http://creativecommons.org/licenses/by/4.0/, as recorded in the arXiv licence metadata for this exact version and shown on the abstract page https://arxiv.org/abs/2301.03486v5. Full licence notice: \texttt{licenses/arxiv-2301.03486-CC-BY-4.0.txt}.\par\smallskip

\noindent\textit{Editorial scope.} Counted unit: the article's lemma lem:diophantine (source0.tex lines 228--230). The article's complete original proof of it is delivered with it (source0.tex lines 232--246, one \texttt{proof} environment of 15 lines). No mathematics is written, repaired or re-derived: the statement and the proof are the article's own lines, in the article's own notation, and no retained line is substituted or re-flowed.  Retained uncounted as the source-local context the proof needs: lines 86--92; lines 180--184.  Nothing was omitted from any retained span, so the package carries no omission record.  No retained line contains a citation, so the wrapper carries no bibliography and no bibliography entry.  No retained line contains a figure environment, an image-inclusion command or drawing code, so no image is delivered and the package declares no asset dependency.  Editorial formatting only, in addition: an AMS \texttt{amsart} preamble that restates the article's own declarations and macros used by the retained lines, namely the environments \texttt{theorem}, \texttt{lemma} on separate counters, together with the standard packages those lines need.  The retained environments print in this document as Theorem 1, Lemma 1 in order of appearance, whereas the article prints the same items with the numbering of its own full document.  The complete native source of the article as distributed in the arXiv e-print for this exact version is bundled unchanged as \texttt{sources/135891/source0.tex}.
\begin{theorem}\label{thm:main}
Let \(p \equiv 1 \pmod{8}\) be prime with \(p^{2} + 1 = 2q\) where \(q\) is prime, and let \(E_{p}: y^{2} = x(x - 1)(x + p^{2})\) be the associated Heronian elliptic curve. Then \(\III(E_{p}/\mathbb{Q})[2]\) is trivial if and only if the Diophantine equation

\[p(x^{2} - py^{2})^{2} = z^{2} + 4x^{2}y^{2}\]

admits a solution in integers with \(z\) odd. Otherwise, \(\III(E_{p}/\mathbb{Q})[2] \cong (\mathbb{Z}/2\mathbb{Z})^{2}\).
\end{theorem}

\begin{align}
p a_{1}^{2} - a_{2}^{2} &= d^{2}, \tag{5}\label{eq:5}\\
p a_{1}^{2} - p a_{3}^{2} &= -p^{2} \cdot d^{2}, \tag{6}\label{eq:6}\\
p a_{3}^{2} - a_{2}^{2} &= 2q \cdot d^{2}. \tag{7}\label{eq:7}
\end{align}

\begin{lemma}\label{lem:diophantine}
\(\alpha = a_{3} + d\sqrt{p} \in \mathbb{Q}(\sqrt{p})\) is a perfect square if and only if the Diophantine equation \(p(x^{2} - py^{2})^{2} = z^{2} + 4x^{2}y^{2}\) is solvable.
\end{lemma}

\begin{proof}
Let us suppose that \(\sqrt{\alpha} \in K\). Then \(\alpha = a_{3} + d\sqrt{p} = (x + y\sqrt{p})^{2}\) for some \(x, y \in \mathbb{Q}\), which implies
\[a_{3} = x^{2} + py^{2}, \qquad a_{1}^{2} = (x^{2} - py^{2})^{2}, \qquad d = 2xy.\]
Substituting into \eqref{eq:7}, we obtain
\[a_{2}^{2} = p(x^{2} - py^{2})^{2} - 4x^{2}y^{2},\]
yielding a solution to the Diophantine equation
\[p(x^{2} - py^{2})^{2} = z^{2} + 4x^{2}y^{2}\]
with \(z\) odd, as \(d\) even implies \(a_{2}\) must be odd.

Conversely, the Diophantine equation \(p(x^{2} - py^{2})^{2} = z^{2} + 4x^{2}y^{2}\) with odd \(z\) is solvable if and only if the homogeneous space corresponding to \((p,1)\) has a rational point
\[\left(\frac{\pm(x^{2} - py^{2})}{2xy}, \frac{\pm z}{2xy}, \frac{x^{2} + py^{2}}{2xy}\right),\]
which is equivalent to the Mordell-Weil rank being at least one. Since the 2-Selmer rank is equal to two, we conclude that
\[r(E_{p}/\mathbb{Q}) = 2,\]
and that the 2-primary component of the Shafarevich-Tate group is trivial. This concludes the proof of Theorem \ref{thm:main}.
\end{proof}

\end{document}
